% GENERATED FILE. Edit canonical lessons, then run npm run generate:books.
% canonical-source-hash: fnv1a64:4696a26725094e57
% canonical-lessons: LA-C434-stabulum, LA-C434-praesepe, LA-C434-faenum, LA-C434-hinnio, LA-C434-mulio

\chapter{Stable, Manger, Hay, To neigh, Muleteer}
\label{ch:la-stable-manger-hay-to-neigh-muleteer}
\hlchaptermodality{434}

\begin{chapteropening}
\textbf{By the end of this chapter:} \emph{I can talk about the stable at an inn: the stable, the manger, hay, horses neighing and the muleteer.}

Complete the last lesson of chapter 434: I can talk about the stable at an inn: the stable, the manger, hay, horses neighing and the muleteer.
\end{chapteropening}

\section[stabulum, stabulī]{stabulum, stabulī — a stable, a stall; a lodging}

\label{lesson:LA-C434-stabulum}



\begin{quote}
Before the new one: say the Latin for a pack saddle, then the Latin for to load.
\end{quote}

\subsection*{You'll want to know: stabulum, stabulī}
\textbf{stabulum, stabulī} — \textquotedblleft{}a stable, a stall; a lodging\textquotedblright{}.

\begin{quote}
Vesperī equī in stabulō quiēscunt.
\end{quote}

\begin{quote}
\emph{In the evening the horses rest in the stable.}
\end{quote}

Second declension, neuter: \emph{stabulum, stabulī n.} Also a cheap roadside lodging where people and beasts stayed. English \emph{stable} comes from it.

\begin{cousinweb}
\emph{stabulum} is built on \emph{stō} 'to stand' with the suffix \emph{-bulum} 'place for' (EtymDB; Lewis \& Short): a standing place. English \emph{stable} comes from it through French.
\end{cousinweb}

\begin{grammarlens}[title={What you've built}]
One more word for telling, explaining and describing at B1.
\end{grammarlens}

\subsection*{Your turn}
\begin{itemize}
\raggedright
  \item \emph{Say it:} \emph{stabulum, stabulī}
  \item \emph{Say it:} \emph{stabulum, stabulī}, once more
  \item \emph{Say it:} a sentence of your own with \emph{stabulum, stabulī}
  \item \emph{Recall:} say the Latin for a pack saddle, then the Latin for to load, then say \emph{stabulum, stabulī} again
\end{itemize}

\begin{tcolorbox}[breakable,colback=teal!4,colframe=teal!35!black,title={Before you move on}]
What does stabulum, stabulī mean? (\textquotedblleft{}A stable, a stall; a lodging\textquotedblright{}.) Say it once more.
\end{tcolorbox}
\section[praesēpe, praesēpis]{praesēpe, praesēpis — a manger, a crib; a stall}

\label{lesson:LA-C434-praesepe}



\begin{quote}
Before the new one: say the Latin for to load, then the Latin for a stable.
\end{quote}

\subsection*{You'll want to know: praesēpe, praesēpis}
\textbf{praesēpe, praesēpis} — \textquotedblleft{}a manger, a crib; a stall\textquotedblright{}.

\begin{quote}
Servus cibum in praesēpe pōnit.
\end{quote}

\begin{quote}
\emph{The slave puts food in the manger.}
\end{quote}

Third declension, neuter: \emph{praesēpe, praesēpis n.} Also spelled \emph{praesaepe}.

\begin{cousinweb}
\emph{praesēpe} is built on \emph{prae-} 'in front' and \emph{saepiō} 'to fence in' (Lewis \& Short, under \emph{praesaepe}): a fenced place in front of the animal. EtymDB lists it without a parent.
\end{cousinweb}

\begin{grammarlens}[title={What you've built}]
One more word for telling, explaining and describing at B1.
\end{grammarlens}

\subsection*{Your turn}
\begin{itemize}
\raggedright
  \item \emph{Say it:} \emph{praesēpe, praesēpis}
  \item \emph{Say it:} \emph{praesēpe, praesēpis}, once more
  \item \emph{Say it:} a sentence of your own with \emph{praesēpe, praesēpis}
  \item \emph{Recall:} say the Latin for to load, then the Latin for a stable, then say \emph{praesēpe, praesēpis} again
\end{itemize}

\begin{tcolorbox}[breakable,colback=teal!4,colframe=teal!35!black,title={Before you move on}]
What does praesēpe, praesēpis mean? (\textquotedblleft{}A manger, a crib; a stall\textquotedblright{}.) Say it once more.
\end{tcolorbox}
\section[faenum, faenī]{faenum, faenī — hay}

\label{lesson:LA-C434-faenum}



\begin{quote}
Before the new one: say the Latin for a stable, then the Latin for a manger.
\end{quote}

\subsection*{You'll want to know: faenum, faenī}
\textbf{faenum, faenī} — \textquotedblleft{}hay\textquotedblright{}.

\begin{quote}
Aestāte agricolae faenum in prātīs colligunt.
\end{quote}

\begin{quote}
\emph{In summer the farmers gather hay in the meadows.}
\end{quote}

Second declension, neuter: \emph{faenum, faenī n.}, also spelled \emph{fēnum}. English \emph{fennel} comes from \emph{faeniculum} 'little hay'.

\begin{cousinweb}
\emph{faenum} is inherited; EtymDB traces it to PIE \emph{d\textsuperscript{h}eh\textsubscript{1}(y)-no-}, and Lewis \& Short tie it to the root of \emph{fēmina} and \emph{fēlīx}. English \emph{fennel} comes from its diminutive \emph{faeniculum} (Lewis \& Short).
\end{cousinweb}

\begin{grammarlens}[title={What you've built}]
One more word for telling, explaining and describing at B1.
\end{grammarlens}

\subsection*{Your turn}
\begin{itemize}
\raggedright
  \item \emph{Say it:} \emph{faenum, faenī}
  \item \emph{Say it:} \emph{faenum, faenī}, once more
  \item \emph{Say it:} a sentence of your own with \emph{faenum, faenī}
  \item \emph{Recall:} say the Latin for a stable, then the Latin for a manger, then say \emph{faenum, faenī} again
\end{itemize}

\begin{tcolorbox}[breakable,colback=teal!4,colframe=teal!35!black,title={Before you move on}]
What does faenum, faenī mean? (\textquotedblleft{}Hay\textquotedblright{}.) Say it once more.
\end{tcolorbox}
\section[hinniō, hinnīre]{hinniō, hinnīre — to neigh, to whinny}

\label{lesson:LA-C434-hinnio}



\begin{quote}
Before the new one: say the Latin for a manger, then the Latin for hay.
\end{quote}

\subsection*{You'll want to know: hinniō, hinnīre}
\textbf{hinniō, hinnīre} — \textquotedblleft{}to neigh, to whinny\textquotedblright{}.

\begin{quote}
In stabulō equī hinniunt, quod cibum exspectant.
\end{quote}

\begin{quote}
\emph{In the stable the horses neigh, because they are waiting for food.}
\end{quote}

Fourth conjugation: \emph{hinniō, hinnīre}; the perfect is rare. A sound word, like English \emph{whinny}.

\begin{cousinweb}
\emph{hinniō} is usually taken as an imitation of the sound a horse makes, like English \emph{whinny}; Lewis \& Short give no origin, and EtymDB lists it without a parent.
\end{cousinweb}

\begin{grammarlens}[title={What you've built}]
One more word for telling, explaining and describing at B1.
\end{grammarlens}

\subsection*{Your turn}
\begin{itemize}
\raggedright
  \item \emph{Say it:} \emph{hinniō, hinnīre}
  \item \emph{Say it:} \emph{hinniō, hinnīre}, once more
  \item \emph{Say it:} a sentence of your own with \emph{hinniō, hinnīre}
  \item \emph{Recall:} say the Latin for a manger, then the Latin for hay, then say \emph{hinniō, hinnīre} again
\end{itemize}

\begin{tcolorbox}[breakable,colback=teal!4,colframe=teal!35!black,title={Before you move on}]
What does hinniō, hinnīre mean? (\textquotedblleft{}To neigh, to whinny\textquotedblright{}.) Say it once more.
\end{tcolorbox}
\section[mūliō, mūliōnis]{mūliō, mūliōnis — a muleteer, a mule driver}

\label{lesson:LA-C434-mulio}



\begin{quote}
Before the new one: say the Latin for hay, then the Latin for to neigh.
\end{quote}

\subsection*{You'll want to know: mūliō, mūliōnis}
\textbf{mūliō, mūliōnis} — \textquotedblleft{}a muleteer, a mule driver\textquotedblright{}.

\begin{quote}
Mūliō mūlōs ad flūmen dūcit et aquam dat.
\end{quote}

\begin{quote}
\emph{The muleteer leads the mules to the river and gives them water.}
\end{quote}

Third declension, masculine: \emph{mūliō, mūliōnis m.} Muleteers hired out mules and carts to travellers.

\begin{cousinweb}
\emph{mūliō} is built on \emph{mūlus} 'mule' (EtymDB; Lewis \& Short) with the suffix \emph{-iō} for a person who works with something.
\end{cousinweb}

\begin{grammarlens}[title={What you've built}]
One more word for telling, explaining and describing at B1.
\end{grammarlens}

\subsection*{Your turn}
\begin{itemize}
\raggedright
  \item \emph{Say it:} \emph{mūliō, mūliōnis}
  \item \emph{Say it:} \emph{mūliō, mūliōnis}, once more
  \item \emph{Say it:} a sentence of your own with \emph{mūliō, mūliōnis}
  \item \emph{Recall:} say the Latin for hay, then the Latin for to neigh, then say \emph{mūliō, mūliōnis} again
\end{itemize}

\begin{tcolorbox}[breakable,colback=teal!4,colframe=teal!35!black,title={Before you move on}]
What does mūliō, mūliōnis mean? (\textquotedblleft{}A muleteer, a mule driver\textquotedblright{}.) Say it once more.
\end{tcolorbox}
